% !TEX root = ../math600-notes.tex

\section{Basics of Groups}

\begin{defn}
    A set $G$ is a \textbf{group} if it's equipped with a binary operation $\cdot :G \times G \to G$ satisfying
    \begin{enumerate}
        \item $(ab)c=a(bc)$
        \item There exists an identity element $e \in G$ such that $ae=ea = a$ for all $a \in G$.
        \item For every $a \in G$ there exists $a^{-1} \in G$ satisfying $aa^{-1}=a^{-1}a=e$.
    \end{enumerate}
\end{defn}

\begin{defn}
    A group is \textbf{abelian} if multiplication is commutative, i.e. $ab=ba$ for all $a,b \in G$. Abeliean groups are often denoted additively, and the identity is denoted by $0$.
\end{defn}

\begin{ex}
    $(\Z, +), (\Q,+)$, $(\R^\times, \cdot)$, $(\C^\times, \cdot)$ are all abelian groups.
\end{ex}

\begin{ex}
    Given a set $A$, we can define $S(A) = \{f:A \to A \mid f$ is bijective$\}$, i.e. the set of all permutations of the elements of $A$. Then, $(S(A), \circ)$ is a group under composition.
\end{ex}

\begin{defn}
    For $A = \{1,\dots,n\}$, $(S(A), \circ)$ is the \textbf{symmetric group} $S_n$
\end{defn}

Note that $S_3$ is non-abelean, and in fact is the smallest non-abelian group

\begin{ex}
    Let $F$ be a field. Then, we have various \textbf{matrix groups}. These include
    \begin{enumerate}
        \item $GL_n(F) = \{ A \in M_n(F) \mid A$ is invertible $\}$. This is called the \textbf{general linear group.}
        \item $SL_n(F) = \{ A \in GL_n(F) \mid \det A = 1\}.$. This is the \textbf{special linear group}
        \item Given any finite dimensional $F$-vector space $V$, we have the group of linear automorphisms $\Aut(V)=\{T \in \Hom(V,V) \mid T$ is invertible$\}$. By choosing a basis $\mathcal B$ of $V$, we can identify $\Aut (V) \cong GL_n(F)$.
        \item When $F=\R$, we have the \textbf{orthogonal group} $O(n) = \{ A \in GL_n(\R) \mid AA^t=I\}$. Moreover, we have the \textbf{special orthogonal group}  $SO(n)= O(n) \cap SL_n(\R)$
        \item When $F = \C$, we have the \textbf{unitary group} $U(n) = \{A \in GL_n(\C) \mid AA^*=I\}$. Moreover, we have the \textbf{special unitary group} $SU(n) = U(n) \cap SL_n(\C)$
    \end{enumerate}
\end{ex}

\begin{prop}
    If $G$ is a group, then
    \hfill
    \begin{enumerate}[(i)]
        \item The identity element of $G$ is unique
        \item Each $a \in G$ has a unique inverse $a^{-1} \in G$
        \item $(a^{-1})^{-1}=a$ for all $a \in G$
        \item $(ab)^{-1}=b^{-1}a^{-1}$ for all $a,b \in G$
    \end{enumerate}
\end{prop}

\begin{proof}

    (ii): Given $a$ with inverses $a',a''$, we have $a' = a'e = a'(a a'')=(a'a)a''=ea''=a''$
\end{proof}

\begin{prop}[Cancellation Laws]
    \begin{enumerate}[(i)]
        \item $ax=ax' \implies x=x'$
        \item $xa=x'a \implies x=x'$
    \end{enumerate}
\end{prop}

\section{Subgroups}

\begin{defn}
    A non-empty subset $H$ of a group $G$ is a \textbf{subgroup} if it's closed under the group operations multiplication and inversion.
\end{defn}

Note that this definition implies $e=aa^{-1} \in H$

\begin{lem}
    Let $H$ be a nonempty subset of a group $G$ which is finite and closed under multiplication. Then, $H$ is a subgroup of $G$.
\end{lem}

\begin{proof}
    Given $a \in H$, we must show that $a^{-1} \in H$. We have $a,a^2,\dots$ are all in $H$, and since $H$ is finite, there exists $r,s$ with $r<s$ such that $a^r=a^s$. Then, $a^{s-r}=e$, so $a(a^{s-r-1})=(a^{s-r-1})a=e$.
\end{proof}

\begin{ex}
    \begin{enumerate}
        \item Every group has the trivial subgroups $\{e\}$ and $G$
        \item Subgroups of $\Z$: Suppose $H$ is a non-zero subgroup of $\Z$. Then, let $n$ be the smallest positive number in $H$. We claim that $H = n\Z=\{ nz \mid z \in \Z\}$.
    \end{enumerate}
    \begin{proof}
        Since $n \in H$ and $H$ is closed under the group operations, the entire subgroup $n\Z \subset H$. Conversely, if $a \in H$, we may divide $a$ by $n$: $a=kn+r$ with $0 \leq r <n$. Then, $r = a-kn\in H$, so by minimality of $n$, we must have $r = 0$. Hence, every $a \in H$ is a multiple of $n$.
    \end{proof}
\end{ex}

\textbf{Application:} Given $a,b \in \Z$ both positive, let $a\Z + b\Z$ be the subgroup of $\Z$ generated by $\{a,b\}$. That is, $a\Z + b\Z = \{ax+by \mid x,y \in \Z\}$. Since this is a nonzero subgroup, we must have that $a\Z + b\Z = d\Z$ for some $d \in \Z^+$. We claim that $d=\gcd(a,b)$.

\begin{proof}
    Indeed, $d \mid a$ since $a=a\cdot 1 + b\cdot 0 \in d\Z$, hence $a = dk$ for some $k \in \Z$. Similarly, $d \mid b$.

    Now, if $c \mid a, c \mid d$, then we must have $c \mid d$, since $d \in a\Z + b\Z$, so $d=ax+by$.
\end{proof}

\begin{defn}
    If $G$ is any group and $a \in G$, we can consider the subgroup \textbf{generated} by $a$. Let $\langle a \rangle = \{a^i : i \in \Z  \}$. If there exists some $a \in G$ such that $G = \langle a \rangle$, we say that $G$ is \textbf{cyclic}.
\end{defn}

\begin{ex}
    For cyclic groups, there are two cases:
    \begin{enumerate}
        \item If all powers $a^i$, $i \in \Z$ are distinct, then $G$ is an infinite cyclic group
        \item If $a^r=a^s$ with $r<s$, then $a^{s-r}=1$. Hence $a^n = 1$ for some smallest positive integer $n$, and so $\langle a \rangle = \{1,a,a^2,\dots,a^{n-1}\}$. We in some cases call this group $C_n$. Note that $C_n$ is realized by the group of $n$'th roots of unity in $\C$, which we denote $\mu_n = \{z \in \C: z^n=1\} = \{e^{2\pi i k/n} \mid k=0,1,\dots,n-1\}$
    \end{enumerate}
\end{ex}

\begin{defn}
    The \textbf{order} of a finite group $G$, denoted $|G|$ or $o(G)$, is the number of elements in $G$. The order $o(a)$ of an element $a \in G$ is the order of the cyclic subgroup $\langle a \rangle$. So, $o(a) = \min \{k \geq 1 \mid a^k=e\}$.
\end{defn}

\begin{defn}
    Let $G$ be a group, $W \subset G$ be a subset. Then, the \textbf{subgroup generated by the subset} $W$, denoted $\langle W \rangle$, is the smallest subgroup of $G$ containing $W$. Concretely, it consists of all products of elements of $W$ and their inverses.
\end{defn}

\begin{defn}
    The \textbf{Klein group} $V$ is the smallest non-cyclic group. It is realized as a subgroup of $GL_2(\R)$ by $V=\left \{ \matr{ \pm1 & 0 \\ 0 & \pm 1} \right \}$. Note that any $2$ non-identity element in $V$ generate $V$.
\end{defn}

\begin{defn}
    The group $GL_n(F)$ has the \textbf{Borel subgroup} $B$ consisting of all upper triangular matrices.
\end{defn}

\section{Cosets and Lagrange's Theorem}

Fun fact: Lagrange invented his theorem before groups were invented

\begin{defn}
    Given $H \leq G$ ("$H$ is a subgroup of $G$"), we define an equivalence relation $x \sim y \iff x^{-1}y \in H$. We say that $x$ \textbf{conguent} to $y$ modulo $H$, and write $x \equiv y \pmod H$

    The equivalence class of $x$ is the (left) \textbf{coset} of $x$, i.e. the set $[x] = \{y \in G: x \sim y\} = \{y \mid x^{-1}y=h \in H\} = \{xh \mid h \in H\}=xH$.
\end{defn}

\begin{notebox*}
    We could have used $xy^{-1}\in H$ in the above definition, and this leads to the right cosets $Hx$
\end{notebox*}

\begin{prop}
    The inversion map $x \mapsto x^{-1}$ gives a $1-1$ correspondence between left and right $H$-cosets, since $H^{-1}=\{h^{-1}\mid h \in H\} = H$, hence $(xH)^{-1}=Hx^{-1}$. We'll use left $H$-cosets. The two notions coincide for abelian groups
\end{prop}

\begin{ex}
    If $G = (\Z, +),H=n\Z$, then $x \equiv y \pmod H \iff x\equiv y \pmod n$. The cosests are $n\Z,1+n\Z,\dots,(n-1)+n\Z$.
\end{ex}

\begin{prop}
    There is a $1-1$ correspondence between any two (left) $H$-cosets. Furthermore, the cosets form a partition of $G$.
\end{prop}

\begin{proof}
    The cosets partition $G$ because they are equivalence classes. Given any two cosets $aH$ and $bH$, the map $f:aH \to bH$ sending $ah \mapsto bh$ is $1-1$ and onto by the cancellation law.
\end{proof}

\begin{thm}[Lagrange's Theorem]
    Let $G$ be a finite group, and $H \leq G$. Then, $|H|$ divides $|G|$
\end{thm}

\begin{proof}
    If $k$ is the number of left $H$-cosets, then $G$ is partitioned into $k$ subsets, each with $|H|$ elements. Hence, $|G| = |H|\cdot k$
\end{proof}

\begin{defn}
    The number of left (or right) $H$-cosets is the \textbf{index} of the subgroup $H$ in $G$, denoted $|G:H|$. Hence, $|G| = |H| \cdot |G:H|$
\end{defn}

\begin{cor}
    \begin{enumerate}
        \item[(a)] If $G$ is a finite group and $a \in G$, then $o(a) \mid o(G)$.
        \item[(b)] If $|G|<\infty$ and $a \in G$, then $a^{|G|}=e$.
        \item[(c)] If $G$ is a finite group with prime order, then $G$ is cyclic.
    \end{enumerate}
\end{cor}

\begin{defn}[The Quotient Set]
    Given any two subsets $A,B$ of a group $G$, we define $AB=\{ab \mid a\in A,b \in B\}$. Let $G/H$ denote the set of all left $H$-cosets of $G$. We want to make this set into a group, with multiplication of cosets.
\end{defn}

Suppose this were the case. Then, for any $g \in G$, $A:=gHg^{-1}H$ must equal another $H$-coset. Clearly, $e \in A$, so as the $H$-cosets partition the group $G$, we must have $A = H$. That is, $gHg^{-1}H=H\implies ghg^{-1}\in H$ for all $g\in G,h \in H$.

\begin{defn}[Normal Subgroups]
    A subgroup $N \leq G$ is a \textbf{normal subgroup} if for each $n \in N$ and $g \in G$, $gng^{-1}\in N$. Equivalently, $gNg^{-1}\subset N$
\end{defn}

\begin{prop}
    The following are equivalent, for $N \leq G$
    \begin{enumerate}
        \item $N$ is a normal subgroup of $G$, denoted $N \unlhd G$
        \item $gNg^{-1}=N$ for all $g \in G$
        \item Every left coset of $N$ is also a right coset of $N$ in $G$.
    \end{enumerate}
\end{prop}

\begin{proof}

    (i) $\implies$ (ii): Assume $gNg^{-1} \subset N$. Then, $N=g(g^{-1}Ng)g^{-1}\subset gNg^{-1}$.

    (ii) $\implies$ (iii): $gNg^{-1}=N \iff gN=Ng$.

    (iii) $\implies$ (i): Given any $g \in G$, we have $gN=Nx$ for some $x \in G$. We must therefore have $Nx=Ng$ since $g\in gN = Nx$ and $g \in Ng$. Hence, we get (i)
\end{proof}

We can therefore define a group structure on $G/N$ when $N$ is normal in $G$, setting $(aN)(bN)=a(Nb)N=a(bN)N=abN$.

\begin{defn}
    The above defines the \textbf{quotient group} $G/N$
\end{defn}

\begin{ex}
    If $G$ is abelian, then every subgroup of $G$ is normal, so we always get a quotient group.
\end{ex}

\begin{ex}[Modular arithmetic]
    Let $G=\Z$, $H=n\Z$. So, $\Z/n\Z=\{\bar 0,\dots,\overline{n-1}\}$, where $\bar x :=x \pmod n$.
\end{ex}

\begin{defn}
    $(\Z / n\Z)^\times$ consists of the multiplicatively invertible elements in $\Z/n\Z \setminus\{0\}$, which are exactly the equivalence classes coprime to $n$. This makes $(\Z/n\Z)^\times$ a multiplicative group with order $\varphi (n)$.
\end{defn}

\begin{notebox*}[Fact]
    Normal subgroups are exactly the kernels of group homomorphisms.
\end{notebox*}

\begin{defn}[Homomorphisms and Isomorphisms]
    A map $\varphi:G \to G'$ between two groups is a \textbf{homomorphism} of groups if it preserves the group operations: $\varphi(ab)=\varphi(a)\varphi(b)$ for all $a,b \in G$. Furthermore, if $\varphi$ is bijective, we say it is an \textbf{isomorphism}.
\end{defn}

\begin{defn}[$\Hom(G,G')$ and $\Aut(G)$]
    We write $\Hom(G,G') :=$ the set of all homomorphisms $G \to G'$. An isomorphism $G\to G$ is called an \textbf{automorphism} of $G$, and these form a group $\Aut(G)$
\end{defn}

\begin{ex}
    For any homomorphism $\varphi :G \to G'$, $\varphi (e)=e'$ and $\varphi(x)^{-1}=\varphi(x^{-1})$.
\end{ex}

\begin{defn}
    The \textbf{kernel} of $\varphi$ is the set of $g \in G$ such that $\varphi (g)=e'$, and the image is $\Img \varphi=\varphi(G)$
\end{defn}

\begin{notebox*}[Check]
    $\ker \varphi$ is always a normal subgroup of $G$.
\end{notebox*}

\begin{thm}[First Isomorphism Theorem]
    For any homomorphism $\varphi:G \to G'$ of groups, $G/\ker \varphi \cong \Img \varphi$, with the map $g\ker \varphi \mapsto \varphi(g)$.
\end{thm}

\begin{proof}
    Same as for vector spaces.
\end{proof}

\begin{ex}
    \begin{enumerate}
        \item For any field $F$, $SL_n(F)=\ker (\det :GL_n(F) \to F^*)$, so we have a normal subgroup of $GL_n(F)$. Furthermore, $GL_n(F) / SL_n(F) \cong (F^*,\cdot)$
        \item If $N \unlhd G$, then we have a projection map $\pi:G \to G/N$ by $g \mapsto gN$, and $N = \ker \varphi$. This helps show how normals subgroups are equivalent to the kernels of homomorphisms
        \item Consider $S^1=\{z \in \C : |z| =1\} \leq \C^\times$, both of which are groups under multiplication. We have $\C^\times \to \R_{>0}$ by $z \mapsto |z|$, and $\ker \varphi = S^1$. So, $\C^\times / S^1 \cong \R_{>0}$
    \end{enumerate}
\end{ex}

\begin{thm}[Universal Property of Quotient Construction]
    % https://q.uiver.app/#q=WzAsMyxbMiwwLCJHIl0sWzQsMiwiRy9OIl0sWzAsMiwiRyciXSxbMCwxLCJcXHBpIl0sWzAsMiwiZiIsMl0sWzEsMiwiXFxleGlzdHMgISBcXHRpbGRlIGYiLDAseyJzdHlsZSI6eyJib2R5Ijp7Im5hbWUiOiJkYXNoZWQifX19XV0=
    Given any group hom $f :G \to G'$ such that $N \subset \ker f$, there exists a unique map $\tilde f:G/N \to G'$ s.t. $\tilde f \circ \pi = f$. That is, the following diagram commutes.
\end{thm}
\[\begin{tikzcd}
	&& G \\
	\\
	{G'} &&&& {G/N}
	\arrow["f"', from=1-3, to=3-1]
	\arrow["\pi", from=1-3, to=3-5]
	\arrow["{\exists ! \tilde f}", dashed, from=3-5, to=3-1]
\end{tikzcd}\]

\section{Products of Groups}

\begin{defn}[External Direct Products]
    Given two groups $G_1,G_2$, we define $G_1 \times G_2$ as a cartesian product of the sets, with the group operation defined component-wise.
\end{defn}

\begin{notebox*}
    $\bar G_1 := G_1 \times \{e_2\}$, and $\bar G_2:= \{e_1\} \times G_2$ are normal subgroups of $G_1 \times G_2$.
\end{notebox*}

\begin{notebox*}[Question]
    Given two subgroups $H$ and $K$ of a group $G$, when is $HK$ a subgroup?
\end{notebox*}

\begin{prop}
    $HK$ is a subgroup of $G \iff HK = KH$
\end{prop}
\begin{proof}
    If $HK \leq G$, then $HK=(HK)^{-1}=K^{-1}H^{-1}=KH$. Now, suppose $HK=KH$. Observe
    $$(h_1k_1)(h_2k_2)=h_1(k_1h_2)k_2=h_1(h_3k_3)k_2=h_4k_4$$
    $$(h_1k_1)^{-1}=k_1^{-1}h_1^{-1}=h'k'.$$
    So, $HK$ is a group
\end{proof}

\begin{cor}
    If $H,K \leq G$ and $H$ or $K$ is normal in $G$, then $HK \leq G$
\end{cor}

\begin{proof}
    Say $H \lhd G$. Then, $HK= \cup Hk=\cup kH=KH$. So, $HK$ is a subgroup by the above proposition
\end{proof}

\begin{defn}[Internal Direct Products]
    Let $H$,$K$ be subgroups of $G$. Then, we say $G$ is the internal direct product of $H$ and $K$ if
    \begin{enumerate}
        \item[(i)] $H$ and $K$ are both normal
        \item[(ii)] $G=HK$
        \item[(iii)] Given any $g \in G$, we can write $g=hk$ for $h \in H$ and $k \in K$ in a unique way
    \end{enumerate}
    This definition generalizes to more than $2$ subgroups.
\end{defn}

\begin{lem}
    Suppose $G=N_1\cdots N_k$ is an internal direct product of its subgroups $N_i \lhd G$. Then, for each $i \neq j$, we have $N_i \cap N_j = \brax{e}$, and for $a \in N_i, b \in N_j$, $ab=ba$.
\end{lem}

\begin{proof}
    Suppose $x \in N_i \cap N_j$. Then, $x=e\cdots e \cdot x_i\cdot e \cdots e=e \cdots e \cdot x_j \cdot e \cdots e$, we by uniqueness of representation, $x=e$. If $a \in N_i$, $b \in N_j$, consider the commutator $[a,b] := aba^{-1}b^{-1}$. We then have $[a,b]=(aba^{-1})b^{-1}=a(ba^{-1}b^{-1})$. Observe $aba^{-1} \in N_j$ and $ba^{-1}b^{-1} \in N_i$ by normality. So, the commutator lies in $N_i \cap N_j = \{e\}$. Thus, $ab=ba$.
\end{proof}

\begin{rmk}
    If $N_1,\dots,N_k$ are normal subgroups of $G$, we need a stronger condition than $N_i \cap N_j = \{e\}$ for $i \neq j$ to ensure that the product is a direct product. Namely, we need $N_i \cap(N_1\cdots N_{i-1}N_{i+1}\cdots N_k)=\{e\}$.
\end{rmk}

\begin{thm}
    Suppose $G$ is the internal direct product $G=N_1\cdots N_k$. Then, $G$ is isomorphic to the external direct product $N_1 \times \cdots \times N_k$.
\end{thm}

\begin{proof}
    The map $\varphi :G\to N_1 \times \cdots \times N_k$ sending $x_1\cdots x_k \mapsto (x_1,x_2,\dots,x_k)$ is well defined by uniqueness. Furthermore, $\varphi$ is clearly bijective. Moreover, $\varphi((x_1\cdots x_k)(y_1\cdots y_k))=\varphi(x_1y_1\cdots x_ky_k)$ by the previous lemma. This maps to
    $$(x_1y_1,\dots,x_ky_k)=(x_1,\dots,x_k)(y_1,\dots,y_k)=\varphi(x_1\cdots x_k)\varphi (y_1\cdots y_k)$$
\end{proof}

Because of this theorem we do not distinguish between internal and external direct products.

\begin{ex}[On Homework]
    If $C_n$ is cyclic of order $n$, show that if $\gcd (m,n)=1$, then $C_m \times C_n \cong C_{mn}$.
\end{ex}

\begin{thm}[Second Isomorphism Theorem]
    If $M$ and $N$ are subgroups of $G$ with $N \lhd G$, then there is a natural isomorphism $MN/N\cong M/M\cap N$
\end{thm}

\begin{proof}
    Let $\varphi:M \to MN/N$ map $m \mapsto mN$. Then, $\varphi$ is clearly surjective, and
    $$\ker \varphi =\{ m \in M \mid m \in N\} = M \cap N.$$
    Thus, the first isomorphism theorem gives $M /M \cap N \cong MN/N$.
\end{proof}

\begin{cor}
    If $M$ and $N$ are as above and are both finite subgroups, then $|MN| = \frac{|M||N|}{|M \cap N|}$
\end{cor}

\section{Inner Automorphisms of a Group}

Let $G = GL_n(F)$. Then, for a fixed $P \in G$, we have the map $A \mapsto PAP^{-1}$, an automorphism. We want to generalize this.

\begin{defn}
    Let $G$ be a group, and $g\in G.$ We define $T_g:G\to G$ by $T_g(x)=gxg^{-1}$. Then, $T_g \in \Aut G$, i.e. it's an isomorphism from $G$ to itself. We call $T_g$ the \textbf{inner automorphism given by $g$}.
\end{defn}

\begin{notebox*}
    $T_g \circ T_h(x)=g(hxh^{-1})g^{-1}=(gh)x(gh)^{-1}=T_{gh}(x)$. So, $T_g \circ T_h = T_{gh}$.
\end{notebox*}

\begin{defn}
    Let $\mathcal I(G)$ be the group of all inner automorphisms of $G$, under composition. We then have a morphism $\psi :G \to \mathcal I(G)$ sending $g \mapsto T_g$.
\end{defn}

What is the kernel of $\psi$?

$\ker \psi = \{g \in G \mid T_g(x) = x\}=\{g \in G\mid gx=xg,x\in G\}=Z(G)$, the \textbf{center} of $G$. Thus,

\begin{prop}
    As $\Img \psi = \mathcal I(G)$, we have $\mathcal I(G) \cong G/Z(G)$.
\end{prop}

\begin{ex}
    For $G = GL_n(F)$, we showed in HW 1 that the center of $Z(G) =\{\lambda I_n:\lambda \in F^\times\}\cong F^\times$. So, $\mathcal I(GL_n(F)) \cong GL_n(F)/F^\times=: PGL_n(F)$. We can think of $GL_n(F)$ as linear automorphisms on $F^n$, and by analogy $PGL_n(F)$ are algebraic automorphisms of $P(F^n)=\mathbb P^{n-1}_F$, which is $n-1$ dimensional projective space.
\end{ex}

\section{Symmetry}

Let $M_n$ be the group of all isometries or rigid motions $\R^n \to \R^n$.

\begin{notebox*}[Recall]
    Any $m \in M_n$ has representation as $m=\tau _v \circ \rho$, where $\tau _v(x) = x+v$ is a translation by $v$, and $\rho \in O(n)$.
\end{notebox*}

\begin{defn}
    Let $F \subset \R^n$ be a subset. A \textbf{symmetry} of $F$ is a rigid motion $m$ with $m(F)=F$. The set of all symmetries of $F$ form a subgroup of $M_n$, which is the group of symmetries of $F$.
\end{defn}

\begin{notebox*}[Goal]
    Find all possible finite groups of symmetry of plane figures.
\end{notebox*}

\begin{thm}
    Let $G$ be a finite subgroup of $M_n$. Then, there exists $x \in \R^n$ fixed by all $g \in G$, i.e. $g(x)=x$ for all $g \in G$.
\end{thm}

\begin{proof}
    Given any point $s \in \R^n$, consider the orbit $\mathcal O_s$ under the action of $G$. That is, $\mathcal O_s:=\{g(s)\mid g \in G\}$. Any $g \in G$ permutes the set $\mathcal O_s$, since $\mathcal O_s$ is a finite set and $g$ is injective on $\R^n$, hence is injective on $\mathcal O_s$.
    \begin{lem}
        Let $S:=\{s_1,\dots,s_n\}$ be a finite set in $\R^m$. Let $p:= \frac 1n(s_1+\cdots +s_n)$ be the \textbf{center of gravity} of $S$. Then, if $s_i'=g(s_i)$ for $g \in M_n$, and $g(p)=p'$, then $p'=\frac 1n(s_1'+\cdots +s_n')$. That is, rigid motions send centers of gravity to centers of gravity.
    \end{lem}
    \begin{proof}
        Since $g = \tau_v \circ \rho$, it suffices to consider the 2 cases:
        \begin{enumerate}
            \item[(i)] $g=\rho$ is linear. Then, $g(p)=g(\frac 1n \sum s_i)=\frac 1n \sum s_i'=p'$.
            \item[(ii)] $g = \tau _v$. Then, $\tau_v(s_i)=s_i+v$, so $\tau _v(p)=p+v=p'$.
        \end{enumerate}
    \end{proof}
    Now, to prove the theorem, notice that the lemma implies any $g \in G$ fixes the center of gravity of $\mathcal O_s$.
\end{proof}

\section{\texorpdfstring{Finite Subgroups of $\mathcal O(2)$}{Finite Subgroups of O(2)}}
By the theorem, any finite subgroup of $M_2$ has a fixed point $o$. By choosing coordinates so that $o$ is the origin, we reduce to $G \leq O(2)$.

\begin{thm}[][label=thm:finite-o2]
    If $G$ is a subgroup of $O(2)$, and $|G| < \infty$, then
    \begin{enumerate}
        \item[(i)] $G= C_n$, the cyclic group of order $n$ generated by the rotation $\rho _\theta$, where $\theta=2\pi /n$
        \item[(ii)] $G=D_n$, the dihedral group of order $2n$, generated by a rotation $\rho_ \theta$ for $\theta = 2\pi /n$, and a reflection $r$ about the line through the origin $0$. We can choose coordinates so that $r$ is a reflection $\matr{1 & 0 \\ 0 & -1}$ about the $x$-axis.
    \end{enumerate}
\end{thm}

\begin{cor}
    If $G$ is a finite subgroup of $M_2$, then for a suitable choice of coordinates, $G$ becomes one of the groups $C_n$ or $D_n$, where $C_n$ is generated by $\rho _{2\pi/n}$, and $D_n$ by $\rho_{2\pi / n}$ and $r$. When $n \geq 3$, $D_n$ is the group of symmetries of an $n$ sided regular polygon.
\end{cor}

\begin{ex}
    \begin{enumerate}
        \item $C_1 = \{1\}, D_1=\{1,r\},C_2=\{1, \rho _ \pi\}, D_2 = \{1, \rho_\pi, r, \rho _\pi \circ r\} \cong V_4$
    \end{enumerate}
\end{ex}

\begin{prop}
    $D_n$ is generated by two elements $x$ and $y$ satisfying the relations $x^n=1$, $y^2=1$, and $yx=x^{-1}y$. The elements of $D_n$ are exactly $\{1,x,x^2,\dots,x^{n-1},xy,x^2y,\dots,x^{n-1}y\}$, so $|D_n| = 2n$.
\end{prop}

\begin{proof}
    We let $x = \rho _{2 \pi /n}$ and $y = r$. Clearly, $x,y$ satisfy the given relations...
\end{proof}

\begin{notebox*}[Exercise]
    If $x=\case{1 \to 3 \\ 2 \to 1 \\ 3 \to 2}$ and $y = \case{1 \to 2 \\ 2 \to 1 \\ 3 \to 3}$ are in $S_3$, check that $x^3=1,y^2=1$, and $yx=x^{-1}y$, hence $D_3 \cong S_3$
\end{notebox*}

\begin{proof}
    Proof of \Cref{thm:finite-o2}. Let $G \leq O(2),|G| < \infty $.

    Case $1$: $G \leq SO(2)$, i.e all elements in $G$ are rotations. If $G = \{1\}$, we're ok, so assume $G$ is non-trivial. Let $\theta = \min \{\alpha >0 \mid \rho _\alpha \in G\}$, which exists as $G$ is finite. Then, $G = \langle \rho _\theta \rangle$. [Indeed, if $\rho _\alpha \in G$, then $\alpha + n \theta + \beta$, where $0 \leq \beta < \theta$. So, $\rho _\beta=\rho _\alpha (\rho _\theta)^{-n}\in G \implies \beta =0$ by definition of $\theta$.] Thus, $G$ is a finite cyclic group $C_n$.

    Case $2$: $G$ contains an element of the form $\rho _\theta \circ r$. Since $\rho _\theta \circ r=\matr{\cos \theta & \sin \theta \\ \sin \theta & -\cos \theta}$ has a square of $1$, it must be a reflection as it is orthogonal. In fact, check that $\rho _\theta \circ r \matr{\cos (\theta /2)\\ \sin (\theta /2)} =\matr{\cos (\theta /2)\\ \sin (\theta /2)}$, so $\rho _\theta r$ is a reflection about the line $\langle \matr{\cos (\theta /2)\\ \sin (\theta /2)} \rangle$. Adjusting coordinates if necessary, we may assume $r \in G$. Let $H:= G \cap SO(2)$, the subgroup of rotations in $G$. By case $1$, we get that $H$ is cyclic, generated by $\rho _\theta$ where $\theta = 2\pi /n$. Hence $G$ contains all products $\rho _\theta^ir^j$ for $0 \leq i < n$, $0 \leq j <2$, so $D_n \subset G$. Now, if $x$ is an arbitrary orientation reversing element of $G$, then $xr \in SO(2) \cap G = H$, so $x \in H_r \subset D_n$. Thus, $G=D_n$.
\end{proof}

\begin{thm}
    Every finite subgroup of $SO(3)$ is one of the following:
    \begin{enumerate}
        \item The cyclic group $C_n$ of rotations by multiples of $2\pi /n$ about a line through the origin.
        \item The dihedral group $D_n$
        \item The tetrahedral group of $12$ rotations carrying a regular tetrahedron to itself. ($A_4$)
        \item The octahedral group of $24$ rotations of a cube or regular octahedron ($S_4$)
        \item The icosahedral group of $60$ rotations of a regular dodecahedron or icosahedron. ($A_5$)
    \end{enumerate}
\end{thm}

In every dimension, the tetrahedron generalizes to the simplex, the cube generalizes to the hypercube, and the regular octahedron generalizes to the convex hull $\langle \pm e_1,\dots, \pm e_n\rangle$.

\begin{rmk}
    The tetrahedron, cube, and octahedron have analogues forms in all dimensions. The only other regular polytopes are
    \begin{enumerate}
        \item Dimension $3$: case (5)
        \item Dimension $4$: (Schläfli, 1842) The $24$-cell, $120$-cell, $600$-cell, bounded by $24$ octahedra, $120$ dodecahedra, and $600$ tetrahedra respectively.
        \item That's it lol
    \end{enumerate}
\end{rmk}

\section{Group Actions}

\begin{defn}[Group Actions]
    Let $G$ be a group and $X$ be a set. We say that $G$ \textbf{acts on} $X$ (on the left) if there exists a map $G \times X \to X$, given by $(g,x) \mapsto g\cdot x$, satisfying
    \begin{enumerate}
        \item[(i)] $e\cdot x=x$ for all $x \in X$
        \item[(ii)] $(g_1g_2)\cdot x=g_1\cdot(g_2\cdot x)$ for all $g_1,g_2\in G,x \in X$.
    \end{enumerate}
    We also say that $X$ is a $G$-set.
\end{defn}

\begin{ex}
    \begin{enumerate}
        \item The group $M_2$ of rigid motions of the plane $\R^2$ acts on the set of points, the set of lines, the set of triangles, etc.
        \item Any group $G$ acts on itself by
        \begin{itemize}
            \item Left multiplication
            \item Conjugation (inner automorphisms)
        \end{itemize}
        \item If $H$ is a subgroup of $G$, then $G$ acts naturally on $G/H$ by $g\cdot(xH)=gxH$.
        \item For any set $X$, $S(X)$ acts on $X$.
        \item The group $G$ of rotational symmetries of the cube acts on the set of (1) faces of the cube; (2) edges of the cube; (3) vertices of the cube.
        \item Suppose $X$ is a connected manifold with universal covering space $p:\tilde X \to X$. Fix a base point $x \in X$. Then the fundamental group $\pi_1(x)$ acts on $p^{-1}(x)$. For example, $p: \R \to S^1$, $\theta \mapsto e^{2\pi i \theta} \dots$
        \item Suppose $K$ is a finite field extension of $\Q$. We have a ring of "integers" $\mathcal O _K$ in $K$ such that $\mathcal O_K \cap \Q = \Z$. Any prime ideal $\mathcal P$ of $\mathcal O_K$ lies above a rational prime $p \in \Z$ if $\mathcal P \cap \Z = p\Z$. For any $p \in \Z$, there exist finitely many prime ideals $P_1,\dots,P_r$ above $p$. If $K/\Q$ is a Galois extension, then the Galois group $G$ acts transitively on the set of these $r$ prime ideals. For example, $K=\Q ( \zeta _n)$, and then $\mathcal O_K=\Z [\zeta _n]$, and the Galois group if $G = (\Z/n\Z)^\times$.
    \end{enumerate}
\end{ex}

\begin{prop}[Permutation Representations]
    Let $G$ be a group and $X$ be a set. There is a $1-1$ correspondence between actions of $G$ on $X$ and group homomorphisms $G \to S(X)$. Such homomorphisms are called \textbf{permutation representations} of $G$
\end{prop}

\begin{proof}
    If $G$ acts on $X$, then for any $g \in G$, define $\eta_g \in S(X)$ by $\eta_g(x):=g\cdot x$. Check that $\eta_g$ is a bijection and the map $g \mapsto \eta_g$ is a group homomorphism $G \to S(X)$. Conversely, given a group homomorphism $\varphi :G \to S(X)$, we define an action of $G$ on $X$ by $g \cdot x = \varphi(g)(x)$. Check that these two processes are inverse to each other.
\end{proof}

\begin{defn}
    A $G$-action is faithful if the corresponding map $G \to S(X)$ is injective. This means that $g\cdot x =x\:\: \forall x \in X \implies g=e$.
\end{defn}

\begin{thm}[Cayley's Theorem]
    Every group $G$ is isomorphic to a subgroup of $S(G)$. In particular, if $|G| = n$, then $G \leq S_n$
\end{thm}

\begin{proof}
    Note that the action of $G$ on itself by left multiplication is faithful, so we get an injection $G \hookrightarrow S(G)$.
\end{proof}

\begin{defn}
    Let $X$ be a $G$-set. Then, for any $x \in X$, the \textbf{orbit} of $x$, denoted $\mathcal O_x$, is the set $\mathcal O_x=\{gx \mid g \in G\}$. $\mathcal O_x$ is itself a $G$-set. Furthermore, the orbits for a $G$-action are equivalence classes for the equivalence relation $x \sim y \iff x=gy$ for some $g \in G$.
\end{defn}

\begin{prop}
    The orbits of $X$ partition the set $X$. Every $x \in X$ lies in exactly one orbit, namely $\mathcal O_x$.
\end{prop}

\begin{defn}
    If $X$ consists of just one orbit, we say that $G$ acts transitively on $X$. So, this occurs if and only if $\forall x,y \in X,\exists g\in G$ such that $y=g \cdot x$.
\end{defn}

\begin{defn}
    If $x,y \in X$ lie in the same $G$-orbit, we often say they are \textbf{conjugates} under the $G$-action.
\end{defn}

\begin{defn}
    For $x \in X$, the \textbf{stabilizer} $G_x$ of $x$ is the subgroup $G_x := \{ g \in G\mid g \cdot x = x\}$
\end{defn}

\begin{thm}
    \begin{enumerate}[(a)]
        \item If $X$ is a transitive $G$-set and $x \in X$, then there is a natural bijection $\varphi:G/G_x \to X$ which maps $gG_x \mapsto g\cdot x$, which is compatible with the $G$-action, that is $\varphi (gC)=g \cdot \varphi(C)$ for every coset $C$. Thus, $G/G_x$ and $X$ are isomorphic as $G$-sets. In particular for an arbitrary $G$-action, we have $G/G_x \cong \mathcal O_x$ for any $x \in X$.
        \item If $X$ is any $G$-set, then $G_{gx}=gG_xg^{-1}$ for all $g\in G,x\in X$. So, the stabilizer subgroups of points belonging to the same $G$-orbit are conjugate.
    \end{enumerate}
\end{thm}

\begin{proof}
    \hfill
    \begin{enumerate}
        \item[(a)]
    \end{enumerate}
    $gG_x = g'G_x \iff g^{-1}g' \in G_x \iff g^{-1}g'(x)=x \iff g \cdot x = g'\cdot x$. This shows $\varphi$ is well defined and injective. If $y \in X$, there exists $g \in G$ such that $y=g\cdot x$ by transitivity, so $y = \varphi(g\cdot G_x)$, thus $\varphi$ is surjective. Finally, it's easy to check that $\varphi$ respects the $G$-action.
    \begin{enumerate}
        \item[(b)]
    \end{enumerate}
    $G_{gx}=\{u \in G \mid ugx = gx\} =\{ u \in G \mid (g^{-1}ug)x=x\}=\{ u \in G\mid g^{-1}ug \in G_x\} = gG_xg^{-1}$
\end{proof}

\begin{ex}
    \begin{enumerate}
        \item Let $M$ be the group of rigid motions of $\R^2$. If $p \in \R^2$ is any point, then $M_p =M_{\tau _p(0)}=\tau _pM_0\tau_p^{-1}=\tau _p O(2) \tau_p^{-1}$, by part (b)
        \item \textbf{(Exercise)} Let $h = \{ z \in \C \mid Im(z) >0\}$ be the upper half plane in $\C$, and $G = SL_2(\R)$. For $g \in \matr{a & b \\ c & d} \in G$, we let $g\cdot z =\frac{az+b}{cz+d}$. Check that this defines a transitive $G$-action on $h$. The stablizer subgroup of $i=\sqrt{-1}$ is $K:=G_i:=\brax{\matr{\cos \theta & -\sin \theta \\ \sin \theta & \cos \theta}\mid \theta \in \R}$. Therefore, $h \equiv SL_2(\R)/K$
    \end{enumerate}
\end{ex}

\begin{cor}
    For any $G$-action on a finite set $X$, we have the counting formulas:
    $$|X|=\sum_{\text{orbits}}|\mathcal O_x|= \sum|G:G_x|,$$
    where the sum is over distinct orbits, and
    $$|\mathcal O_x| = |G| / |G_x| \iff |G| = |\mathcal O_x| \cdot |G_x|,\forall x \in X$$
\end{cor}

\begin{ex}
    \begin{enumerate}
        \item A group is \textbf{simple} if its only normal subgroups are $\{e\}$ and itself. If $H$ is any proper subgroup of a simple group $G$, we get an injection $G \hookrightarrow S_k$, where $k = |G:H|$ (Just consider the $G$-action on the coset space $G/H \ \iff$ a map $G \to S(G/H)$)
        \item \textbf{Check:} If $H$ and $K$ are subgroups of $G$, then the $G$-sets $G/H$ and $G/K$ are isomorphic if and only if $H$ and $K$ are conjugate subgroups of $G$.
        \item Suppose $G$ is a finite group and $H$ and $K$ are two subgroups of $G$. We compute $|HK|$. Consider $X := G/K$ and consider the action of $H$ on $X$ by left multiplication. Then, the orbit of $K$ in $G/K$ under the $H$-action is $HK$. Therefore, $|HK| = |K| \cdot$(\# of $H$-cosets that lie in this orbit). The stabilizer $H_K$ of the coset $K$ is $H \cap K$, hence the orbit in question includes $|H:H \cap K|$ cosets of $K$. Thus, $|HK| = \frac{|H| \cdot |K|}{|H \cap K|}$
    \end{enumerate}
\end{ex}

\section{\texorpdfstring{Conjugacy Classes in $G$; The Class Equation}{Conjugacy Classes in G; The Class Equation}}
Recall the conjugation action of a group $G$ on itself is defined by $g \cdot x = gxg^{-1}$.
\begin{defn}
    The stabilizer of an element $x \in G$ under the action of conjugation is called the \textbf{centralizer} of $x$, denoted $Z(x)$. That is, $Z(x)=\{ g \in G \mid gx=xg\}$.
\end{defn}

\begin{defn}
    The orbit $\mathcal O _x$ of $x$ under the $G$-action is the \textbf{conjugacy class} $C_x=\{ y \in G \mid y= gxg^{-1}$ for some $g \in G \}$
\end{defn}

\begin{prop}[The Class Equation]
    By the counting formula, $|G| = |Z(x)| \cdot |C_x|$, so $|C_x| = |G:Z(x)|$

    We also have the \textbf{class equation}: let $x_1,\dots,x_n$ be all the representatives of distinct conjugacy classes, then
    $$|G| = \sum_{i=1}^n|C_{x_i}|=\sum_{i=1}^n|G:Z(x_i)|.$$
    Furthermore, any element in the center of $G$ only has itself in the orbit. Therefore, we can let $g_1,\dots,g_m$ be all representatives of non-central conjugacy classes, and then
    $$|G|=|Z(G)|+\sum_{i=1}^m|G:Z(g_i)|$$
    Note by assumption that $|G:Z(g_i)|\geq 2$.
\end{prop}

\begin{rmk}
    All the integers on the right hand side of the class equation divide $|G|$, and $|Z(G)| \geq 1$.
\end{rmk}

\begin{ex}
    $G= D_3$ has conjugacy classes $\{1\},\{r,r^2\},\{s,sr,sr^2\}$. The class equation is $6 = 1+2+3$.
\end{ex}

\begin{ex}
    If $X$ is the set of nilpotent matrices in $M_n(F)$ and $G =GL_n(F)$ acting on $X$ by conjugation, then the $G$-orbits are given by Jordan canonical form. We'll prove this later.
\end{ex}

\begin{defn}
    A group $G$ is a \textbf{$p$-group} if $|G|=p^k$, $p$ prime.
\end{defn}

\begin{prop}
    The center $Z(G)$ of any $p$-group $G$ is non-trivial.
\end{prop}

\begin{proof}
    If $Z(G)=\{1\}$, then the class equation for $G$ reads $p^k=|G|=1+\sum |G:Z(g_i)|$, but $|G:Z(g_i)|=p^{k_i}$ for some $k_i>1$, which is a contradiction.
\end{proof}

\begin{cor}
    Every group of order $p^2$ for a prime $p$ is abelian.
\end{cor}

\begin{proof}
    Suppose $|G|=p^2$. If $Z(G)=G$, then $G$ is abelian. Otherwise, we have $|Z(G)|=p$. Let $x \in G$, and we'll show $Z(x)=G$. If $x \in Z(G)$, we're ok, and otherwise we have $1 \subsetneq Z(G) \subsetneq Z(x)$, which implies $Z(x)=G$.
\end{proof}

\begin{cor}
    If $|G|=p^2$, then
    \begin{enumerate}
        \item[(i)] $G=C_{p^2}$ is cyclic, or
        \item[(ii)] $G = C_p \times C_p$
    \end{enumerate}
\end{cor}

\begin{proof}
    We know that $G$ is abelian. If $G$ contains an element of order $p^2$, then $G$ is cyclic. Otherwise, the order of any non-identity element $x$ of $G$ is $p$, so pick $x \neq 1$, set $H_1 = \langle x \rangle$. Pick $y \in G \setminus H_1$, and $H_2 = \langle y \rangle$. Then, $H_1 \unlhd G$ and $H_2 \unlhd G$ by $G$ being abelian. Furthermore, $H_1 \cap H_2 = \{1\}$, and since $H_1 \subsetneq H_1H_2$, we have $H_1H_2=G$. So, $G \cong H_1 \times H_2=C_p \times C_p$.
\end{proof}

\begin{rmk}
    Groups of order $p^3$ are more complicated
\end{rmk}
There is $C_{p^3},C_{p^2}\times C_p,C_p\times C_p \times C_p$, which are distinct. However, there are also some non-abelian groups. In fact, for any prime $p$, there are exactly two non-isomorphic non-abelian groups of order $p^3$.

\begin{ex}[$p=2$]
    When $p=2$, there are $2$ non-abelian groups of order $8$. They are the dihedral group $D_4$, and the \textbf{quaternion group} $Q=\{\pm 1, \pm i, \pm j, \pm k\mid(-1)x=-x,i^2=j^2=k^2=ijk=-1 \}$. $Q$ has only one element of order $2$, so $Q \not \cong D_4$.

    Here is a presentation of $Q$ using matrices: Let $a = \matr{0 & \sqrt {-1} \\\ \sqrt{-1} & 0},b=\matr{0 & 1 \\ -1 & 0}$. Then, $a^4=1,a^2=b^2$, and $ba=a^3b$. We get that $Q = \langle a,b \rangle$. Prove this classification as an \textbf{exercise}.
\end{ex}

If $p$ is odd, then the non-abelian groups of order $p^3$ are
$$Heis(\Z/p\Z) =\{ \matr{1 & a & b \\ 0 & 1 & c \\ 0 & 0 & 1}\mid a,b,c\in \Z/p\Z\}$$
$$G_p=\{ \matr{a & b \\ 0 & 1} \mid a,b \in \Z/p^2\Z,a \equiv 1 \pmod p \}$$
Note that when $p=2$, the Heisenberg group and $G_2$ are both isomorphic to $D_4$.

\section{\texorpdfstring{Conjugacy Action of $G$ on Subgroups}{Conjugacy Action of G on Subgroups}}
$G$ acts on the set $X = \{$ subgroups of $G\}$ by conjugation: $g\cdot H:= gHg^{-1}$.

\begin{defn}
    The stabilizer of a subgroup $H$ under this action is $G_H=\{g\in G \mid gHg^{-1}=H\}$, which we call the \textbf{normalizer} $N_G(H)$.
\end{defn}

\begin{rmk}
    Clearly, $H \lhd N(H)$, and one can check that $N(H)$ is the largest subgroup of $G$ which contains $H$ as a normal subgroup.
\end{rmk}

\begin{cor}
    The number of conjugate subgroups of $H$ equals the index of the normalizer $N_G(H)$ in $G$.
\end{cor}

\begin{notebox*}[Notation]
    If $G$ acts on a set $X$, then let $X^G$ denote the set of fixed points $\{x \in X \mid gx=x\}$. Then, $|X|=|X^G|+\sum_{|\mathcal O_i|>1}|\mathcal O_i|=|X^G|+\sum_{G_{x_i} \neq G} |G:G_{x_i}|$
\end{notebox*}

\section{The Sylow Theorems}
The Sylow theorems give a partial converse to Lagrange's theorem.

\begin{defn}
    Suppose $|G|=p^rm$, where $\gcd(p,m)=1$. A \textbf{$p$-Sylow subgroup} of $G$ is a subgroup of order $p^r$. A subgroup of $G$ is called a \textbf{Sylow subgroup} if it is a $p$-Sylow subgroup for some $p$ dividing $|G|$
\end{defn}

\begin{thm}[First Sylow Theorem]
    For every prime $p$ dividing $|G|$, $G$ has a $p$-Sylow subgroup
\end{thm}

\begin{proof}
    Let $|G|=p^rm$, $r\geq 1$ and $p \nmid m$. Let $X$ be the collection of all subsets of $G$ of cardinality $p^r$. That is, $X:= \{A \subset G : |A| = p^r\}$. $G$ acts on $X$ by left multiplication $A \mapsto g\cdot A$.

    \begin{lem}
        The cardinality $|X|=\binom{p^r m}{p^r} \not \equiv 0 \pmod p$.
    \end{lem}
    \begin{proof}
        $$\binom{p^r m}{p^r}=\frac{p^rm(p^rm-1)\cdots (p^rm-p^r+1)}{p^r\cdot (p^r-1)\cdots 1}$$
        Note that for each integer $k$ with $0 \leq k < p^r$, $p^i\mid (p^rm-k) \iff p^i \mid (p^r-k)$. (Indeed, if we write $k=p^jl,j<r,p \nmid l$, then $p^rm-k = p^rm-p^jl = p^j(p^{r-j}m-l)$). This implies that all powers of $p$ in the numerator are exactly canceled by the powers of $p$ in the denominator.
    \end{proof}
    Now, write $X$ as a union of $G$-orbits $\mathcal O$. Then, $|X|=\sum|\mathcal O|$. Since $p \nmid |X|$, there exists some orbit $\mathcal O$ with $|\mathcal O| \not \equiv 0 \pmod p$. Let $A \in \mathcal O$ be a subset of $\mathcal O$ with $1 \in A$. (Such an $A$ lies in every orbit). Consider $G_A$. We have $G_A \subset G_A \cdot A=A \implies |G_A| \leq |A| =p^r$. On the other hand, $|G|=|\mathcal O|\cdot |G_A|$. So, as $p \nmid \mathcal |O|$, we must have $p^r \mid |G_A|$, showing $|G_A| \geq p^r$. Hence, $|G_A|= p^r$, so $G_A$ is a $p$-Sylow subgroup.
\end{proof}

\begin{rmk}
    A variant of this argument shows that $G$ has a subgroup of order $p^i$ for each $i\in \{1,\dots,r\}$.
\end{rmk}

\begin{cor}[Cauchy's Theorem]
    If $p$ divides $|G|$, then $G$ contains an element of order $p$.
\end{cor}
\begin{proof}
    Let $H$ be a $p$-Sylow subgroup of $G$, and $x \in H$, $x \neq 1$. Then, $|x|=p^k$ for some $k \geq 1$. Hence, $(x^{p^{k-1}})^p=x^{p^k}=1$, and by minimality of order $x^{p^{k-1}}\neq 1$, showing $|x^{p^{k-1}}|=p$.
\end{proof}

\begin{thm}[Second Sylow Theorem]
    Any two $p$-Sylow subgroups of $G$ are conjugate.
\end{thm}

\begin{proof}
    We'll show something stronger: Let $H \leq G$ be any $p$-subgroup of $G$ and $P$ be a $p$-Sylow subgroup. We claim that there exists $x \in G$ such that $H \subset xPx^{-1}$.

    Consider $G/P$, the set of all $P$ cosets $xP$ in $G$, and the left $H$-action on $G/P$ by $h \cdot xP=hxP$ for all $xP \in G/P$. By the counting formula, recalling $|G|=m\cdot p^r=m|P|$, we get that
    $$m = |G/P|=|(G/P)^H|+\sum_{H_i \neq H} |H:H_i|.$$
    Since $H$ is a $p$-group all $|H:H_i|$ are non trivial powers of $p$. So, modulo $p$, this implies $0 \not \equiv m \equiv |(G/P)^H| \pmod p$. We deduce in particular that $(G/P)^H$, the set of $H$-fixed points, is not empty. So, there exists $xP$ fixed by all $h \in H$. That is,
    $$hxP=xP\implies x^{-1}hx \in P \implies h \in xPx^{-1} \implies H \subset xPx^{-1}$$
\end{proof}

\begin{cor}
    \begin{enumerate}[(i)]
        \item All Sylow $p$-subgroups of $G$ are isomorphic
        \item Let $P$ be a $p$-Sylow subgroup. Then, $P \lhd G \iff P$ is the unique $p$-Sylow subgroup.
        \item The number of $p$-Sylow subgroups is $|G:N_G(P)|$, where $P$ is a $p$-Sylow subgroup, and note that $|G:N_G(P)|$ divides $|G|$.
        \item Any $p$-group in $G$ is contained in a $p$-Sylow subgroup.
    \end{enumerate}
\end{cor}

\begin{cor}
    Let $p$ be a prime, and let $|G|_p$ denote the highest power of $p$ that divides $|G|$. Then, the number of $p$-Sylow subgroups of $G$ divides $|G|/|G|_p$.
\end{cor}
\begin{proof}
    We know that the number of $p$-Sylow subgroups of $G$ is $|G:N_G(P)|$, where $P$ is any one of them. Moreover, $|G:N_G(P)| \cdot |N_G(P):P|=|G:P|=|G|/|P|=|G|/|G|_p$. Thus, $|G:N_G(P)|$ divides $|G|/|G|_p$ as desired.
\end{proof}

\begin{cor}
    Any element of order $p$ in $G$ is contained in some $p$-Sylow subgroup.
\end{cor}

\begin{thm}[Third Sylow Theorem]
    The number of $p$-Sylow subgroups is $\equiv 1 \pmod p$.
\end{thm}

\begin{proof}
    Let $X$ be the set of all $p$-Sylow subgroups, and let $P$ be one of them. Then, $P$ acts on $X$ by conjugation, and for this action, we have $|X|\equiv |X^P| \pmod p$ by the counting formula (recall that all the other terms are divisible by $p$). We claim that $|X^P|=1$, in other words $X^P = \{P\}$.

    Indeed, suppose $Q \in X^P$, so $xQx^{-1}=Q$ for all $x \in P$. That is, $P \subset N_G(Q)$. Now, $Q$ is a $p$-Sylow subgroup of $N_G(Q)$, and $Q \unlhd N_G(Q)$. Hence, $P=Q$.
\end{proof}

Note that the number of $p$-Sylow subgroups must be $1$ mod $p$ and must divide the order of $|G|$, which is a fairly strong condition.

\begin{ex}
    Every group of order $15$ is cyclic:
    \begin{proof}
        The number of $3$-Sylow subgroups divides $15/3 = 5$ and is $\equiv 1 \pmod 3$, so there exists a unique $3$-Sylow subgroup $H$. Similarly, the number of $5$-Sylow subgroups divides $15/5=3$ and is $\equiv 1 \pmod 5$, so there's a unique $5$-Sylow subgroup $K$. The uniqueness implies $H,K \unlhd G$, and clearly $H\cap K = 1,HK=G$, so $G \cong H \times K \cong C_3 \times C_5 \cong C_{15}$.
    \end{proof}

    This generalizes to groups of order $pq$ with $q \not \equiv 1 \pmod p$.
\end{ex}

\begin{ex}
    Any finite abelian group is the direct product of its Sylow subgroups.
    \begin{proof}
        If $P_1,\dots,P_k$ are the distinct Sylow subgroups of an abelian group $A$, then $P_1,\dots,P_k \unlhd A$ by $A$ being abelian. Moreover, $P_i \cap (P_1\cdots \hat P_i \cdots P_k) =1$ by order considerations, for all $i=1,\dots,k$. Furthermore, $P_1P_2\cdots P_k=A$, also by cardinality considerations. So, $G=P_1\times \cdots \times P_k$.
    \end{proof}
\end{ex}

Later: we'll classify all finite abelian groups.

\section{The Symmetric Group}

Recall that $S_n = S(\{1,\dots,n\})$, with order $n!$.

\begin{defn}
    A permutation $\sigma \in S_n$ is a \textbf{cycle of length} $r$, or an $r$\textbf{-cycle} if there exists distinct integers $a_1,\dots,a_r \in [1,n]$ such that $\sigma(a_i)=a_{i+1}$ for $i < r$, and $\sigma(a_r) = a_1$, while $\sigma(x)=x$ for all $x \not \in \{a_1,\dots,a_r\}$. We write $\sigma =(a_1 \cdots a_r)$, and we note that any $r$-cycle can be written like this in $r$ different ways. We say that two cycles are \textbf{disjoint} if no number is moved by both of them.
\end{defn}

\begin{prop}
    Any two disjoint cycles commute. Moreover, any $\sigma \in S_n$ has a (unique) decomposition as a product of disjoint cycles.
\end{prop}

\begin{proof}
    (Sketch): Given any $w \in S_n$. we consider the subgroup $\angles{w}$ generated by $w$. The cycles in the decomposition of $w$ correspond to the orbits of the $\angles{w}$-action on the set $\{1,2,\dots,n\}$.
\end{proof}

\begin{ex}
    $\sigma = 46835217 \in S_8$ has cycle decomposition $(14387)(26)$.
\end{ex}

\begin{defn}
    A \textbf{transposition} in $S_n$ is any $2$-cycle (ij).
\end{defn}

\begin{prop}
    Every $w \in S_n$ is a product of (not necessarily disjoint) transpositions (in a non-unique way)
\end{prop}

\begin{proof}
    Any $r$-cycle $(a_1\cdots a_r)$ satisfies $(a_1\cdots a_r)=(a_1a_r)(a_1a_{r-1})\cdots (a_1a_2)$.
\end{proof}

\begin{defn}
    A permutation is \textbf{even} if it can be written as a product of an even number of transpositions. Otherwise, it is \textbf{odd}.
\end{defn}
We'll prove that this notion is well defined.

\section{\texorpdfstring{Conjugacy in $S_n$}{Conjugacy in Sn}}

\begin{prop}
    Let $\sigma,\tau \in S_n$, and suppose that $\sigma$ has cycle decomposition
    $$\sigma = (a_1a_2\cdots a_{k_1})(b_1b_2\cdots b_{k_2})\cdots$$
    Then, $\tau \sigma \tau ^{-1}=(\tau (a_1)\cdots (\tau(a_{k_1}))(\tau (b_1) \cdots (\tau (b_{k_2}))$, hence two elements of $S_n$ are conjugate if and only if their cycle decompositions have the same orders (i.e. $\sigma$ and $\tau$ have the same cycle type). The number of conjugacy classes in $S_n$ equals the number of integer partitions of $n$.
\end{prop}

\begin{proof}
    It suffices to check that $\tau (ij)\tau ^{-1}=(\tau (i)\tau (j))$ for any transposition $(ij)$ (easy). Alternatively, observe that if $\sigma (a)=b$, then $\tau\sigma \tau ^{-1}(\tau (a))=\tau (b)$. Hence, if the pair "ab" appears in the cycle decomposition of $\sigma$, then the pair $\tau(a)\tau (b)$ appears in the cycle decomposition of $\tau \sigma \tau ^{-1}$.
\end{proof}

\begin{cor}
    If $n \geq 3$, then $Z(S_n)$ is trivial.
\end{cor}

\begin{ex}[Class Equation for $S_4$]
    The conjugacy classes are $1$, the $6$ transpositions, the $3$ transposition products, the $8$ three cycles, and the $6$ four cycles. So, $|S_4| =24=1+3+6+6+8$
\end{ex}

Back to the proof that "even" and "odd" are well defined.

\begin{proof}
    We consider the Vandermonde product $\Delta (x) = \prod _{1 \leq i < j \leq n}(x_i-x_j)$, where $x_1,\dots,x_n$ are variables. If $f(x_1,\dots,x_n)$ is any polynomial in $x_1,\dots,x_n$, the symmetric group $S_n$ acts on $f$ by permuting the variables. That is, $\sigma \cdot f(x_1,\dots,x_n)=f(x_{\sigma(1)},\dots,x_{\sigma (n)})$. So, observe $\sigma \Delta = \prod_{1 \leq i < j \leq n}(x_{\sigma(i)}-x_{\sigma (j)})= \pm \Delta$. Now, define $\varphi : S_n \to \{ \pm 1\}$ by $\varphi(\sigma) = \case{+1 & \sigma \Delta = +\Delta \\ -1 & \sigma \Delta = -\Delta}$. We claim $\varphi$ is a group homomorphism. Indeed, $(\sigma \tau)(\Delta(x))=\prod_{ 1\leq i < j \leq n}(x_{\sigma \tau (i)}-x_{\sigma \tau (j)}) = \sigma \left ( \prod_{1 \leq i < j \leq n}(x_{\tau (i)} - x_{\tau (j)}) \right )$. It follows that $\varphi (\sigma \tau) = \varphi (\sigma) \varphi (\tau)$.

    Next, we claim that for any transposition $\tau =(ab)$, $\tau (\Delta)=-\Delta$. This is obvious when $\tau = (12)$, and in general, if $\tau = (ab)$, let $w \in S_n$ be any element such that $w(1)=a,w(2)=b$. Then, $w(12)w^{-1}=(ab)$. Since $\varphi$ is a homomorphism, it follows that
    $$\varphi (ab)=\varphi (w)\varphi (12) \varphi (w)^{-1}=\varphi (12)=-1$$
    Finally, if $\sigma \in S_n$ is any permutation, and $\sigma = \tau_1 \dots \tau _k$, where each $\tau _i$ is a transposition, then $\varphi (\sigma)=\varphi(\tau_1)\cdots \varphi(\tau _k)=(-1)^k$, hence the parity of $k$ depends only on $\sigma$.
\end{proof}

\begin{defn}
    The \textbf{alternating group} $A_n = \ker \varphi$, i.e. the group of all even permutations in $S_n$.
\end{defn}

\begin{prop}
    $|S_n : A_n|=2$, and $A_n \unlhd S_n$.
\end{prop}

\begin{thm}[Simplicity of $A_5$ (Galois)]
    $A_5$ is a simple group.
\end{thm}

\begin{proof}
    We have $|A_5|=60=2^2\cdot 3\cdot 5$. The number of Sylow $5$-subgroups of $A_5$ is a divisor of $60/5=12$ and is $\equiv 1 \pmod 5$. Since any $5$-cycle is $A_5$ generates a Sylow $5$-subgroup, we see that the number of Sylow $5$-subgroups if $6$. Hence, the number of elements of order $5$ in $A_5$ is $6\cdot (5-1)=24$. These are the $5$-cycles $(1xyzw)$ where $xyzw$ is a permutation of $2345$.

    Similarly, $A_5$ could have $1,4,$ or $10$ Sylow $3$-subgroups, and by inspection we see the number is $10$. These $10$ subgroups given $10\cdot(3-1)=20$ elements of order $3$ in $A_5$.

    Choose $i \in \{1,\dots,5\}$, let $\{a,b,c,d\}=\{1,2,3,4,5\}\setminus\{i\}$, and define
    $$V_i:=\{1,(ab)(cd),(ac)(bd),(ad)(bc)\} \leq A_5.$$
    Each $V_i$ is a Sylow $2$-subgroup of $A_5$, and $V_i \cap V_j = \{1\}$ for $i \neq j$, and by inspection $\sigma V_i\sigma ^{-1}=V_{\sigma(i)}$ for all $\sigma \in A_5$. Therefore, $V_1,\dots,V_5$ are the only Sylow $2$-subgroups of $A_5$. It also follows that $A_5$ has $5 \cdot (4-1)=15$ elements of order $2$. One checks (exercise) that any $2$ such elements of $A_5$ are conjugate to each other.

    Now, let $N \unlhd A_5$ be nontrivial. If $N$ contains an element $x$ of order $5$, then it contains the Sylow $5$-subgroup generated by $x$, hence (since $N$ is normal) all of the Sylow $5$-subgroups of $A_5$ are in $N$. That is, $N$ either contains all $24$ elements of order $5$ or none of them. Similarly, $N$ contains all $20$ elements of order $3$ or none of them, and $N$ contains all $15$ elements of order $2$ or none of them. But then the cardinality is $1+|X|$, where $X$ is a nonempty subset of $\{15,20,24\}$. Also, $|N|$ divides $60$, so we deduce that $|N|=60$, thus $N=A_5$.
\end{proof}

\section{Finite Simple Groups}
\begin{lem}
    $A_5$ is the only nontrivial proper normal subgroup of $S_5$.
\end{lem}

\begin{proof}
    Suppose $H \unlhd S_5, \{1\} \neq H \neq S_5$. If $H \neq A_5$, then $H \not \subset A_5$ because $A_5$ is simple, so $HA_5 = S_5$. The second isomorphism theorem gives $HA_5/H \cong A_5/(H \cap A_5)$. Therefore, $H \cap A_5$ is a proper normal subgroup of $A_5$, hence $H\cap A_5 = \{1\}$. It follows that $|HA_5|/|H|=|A_5|$, hence $|H|=2$. But, no subgroup of order $2$ in $S_5$ can be normal.
\end{proof}

\begin{thm}
    Any simple group of order $60$ is isomorphic to $A_5$.
\end{thm}

\begin{proof}
    Let $G$ be a simple group of order $60$. If $G$ has a subgroup $H$ with index $|G:H|=n$, then the left action of $G$ on $G/H$ gives a homomorphism from $G$ to $S(G/H) \cong S_n$. Since $G$ is simple, $\ker \varphi =1$ or $G$, but clearly $\ker \varphi$ can't be $G$. So, $\varphi$ is injective. Thus, if $n \leq 4$, then $|S_n| \leq 24 <60$, so $G$ couldn't be a subgroup of $S_n$. We deduce that $H=G$ or $|G:H| \geq 5$.

    We now claim that $G$ has a subgroup of index $5$, and hence $G$ is isomorphic to a subgroup of $S_5$. This claim will prove the theorem, since $G$ would have index $2$ in $S_5$, hence $G \unlhd S_5$, and we use the lemma.

    Suppose $G$ does not have a subgroup of index $5$. The number of Sylow $2$-subgroups of $G$ divides $60/4=15$. Since this number equals $|G:N(P)|$ for any Sylow $2$-subgroup $P$, by our hypothesis, and the previous analysis, we deduce that the number is exactly $15$. Let $H_1$ and $H_2$ be distinct $2$-Sylow subgroups of $G$, and suppose there exists $1 \neq t \in H_1 \cap H_2$. Then, $|Z(t)|>4$, because $H_1$ and $H_2$ are distinct and abelian, both contained in $Z(t)$. Also $4$ divides $|Z(t)|$ since $H_1 \subset Z(t)$. It follows that $|G:Z(t)| \leq 5$. Previous analysis and our assumption shows that $Z(t)=G$, but then $t \in Z(G)$, so $G$ has a nontrivial center, contradicting the fact that $G$ is simple. Therefore, no two of the $15$ $2$-Sylow subgroups of $G$ have a non-trivial element in common. Hence, $G$ has $15\cdot (4-1)=45$ elements of order $2$ or $4$. But, $G$ is simple, so $G$ must have more than $1$ $5$-Sylow subgroup, this $G$ has $6$ $5$-Sylow subgroups, and hence $24$ elements of order $5$. Bad!
\end{proof}

\begin{rmk}
    It is a relatively easy (and tedious) exercise to show that the only simple groups of order $<60$ are the cyclic groups of prime order. Hence, $A_5$ is the smallest non-abelian simple group. The non-abelian finite simple group of next smallest order is $PSL_2(\F_7)$, which has order $168$. (Note: $A_5 \cong PSL_2(\F_5)\cong PSL_2(\F_4)$)
\end{rmk}

\begin{thm}[Classification of Finite Simple Groups, 1984, 15000 pages]
    \begin{enumerate}
        \item $(\Z/p\Z,+)$, $p$ prime
        \item $A_n,n \geq 5$
        \item Finite groups of Lie type, such as $PSL_2(\Z/p\Z)$, $p\geq 5$ prime.
        \item $26$ others, known as \textbf{sporadic groups}. The smallest is the Matthieu group $M_{11}$ of order $7920$, and the largest is the monster group of order $\sim 10^{53}$
    \end{enumerate}
\end{thm}

\begin{defn}
    A \textbf{composition} series for a group $G$ is a normal series $1 = G_s \lhd\cdots \lhd G_1 \lhd G_0=G$ such that each $G_{i+1}$ is a maximal proper normal subgroup of $G_i$, i.e. $G_{i}/G_{i+1}$ is a nontrivial simple group for all $i \geq 1$. These quotient groups are the \textbf{composition factors} of the series.
\end{defn}

\begin{rmk}
    Clearly, finite groups have a composition series.
\end{rmk}

\begin{thm}[Jordan-Hölder Theorem]
    Let $G$ be a finite group and $1=G_s \lhd \cdots \lhd G_1 \lhd G_0 = G$ and $1=H_t \lhd \cdots \lhd H_1 \lhd H_0 = G$ be two composition series for $G$. Then, $s=t$, and there exists a permutation $i \to j$ of $[1,s]$ such that $G_i/G_{i+1} \cong H_j/H_{j+1}$ for all $1 \leq i \leq s$. Thus, the composition factors $G$ are determined up to isomorphism.
\end{thm}

\begin{proof}
    We prove this by induction on $|G|$. This is trivial for groups of small order. We have to cases:
    \begin{enumerate}
        \item $G_1 =H_1$: The claim is clear by induction
        \item $G_1 \neq H_1$:
    \end{enumerate}
    Since $G_1,H_1 \lhd G$, we have $G_1H_1 \lhd G$ As $G_1  \subsetneq G_1H_1$, we have $G_1H_1 =G$ by maximality of $G_1$ as a normal subgroup. So, by the 2nd isomorphism theorem, we get $G/G_1 \cong H_1/(G_1 \cap H_1)$ and similarly $G/H_1 \cong G_1 /(G_1 \cap H_1)$. Let $K_2 := G_1 \cap H_1$. Then, $K_2$ is a maximal normal subgroup of both $G_1$ and $H_1$. Let $1=K_u \lhd \cdots \lhd K_3 \lhd K_2$ be a composition series for $K_2$. We then have the following $4$ composition series:
    \begin{enumerate}[(i)]
        \item $G_s \lhd \cdots \lhd G_2 \lhd G_1 \lhd G_0=G$
        \item $K_u \lhd \cdots \lhd K_2 \lhd G_1 \lhd G_0 = G$
        \item $K_u \lhd \cdots \lhd K_2 \lhd H_1 \lhd H_0 =G$
        \item $H_t \lhd \cdots \lhd H_2 \lhd H_1 \lhd H_0 = G$
    \end{enumerate}
    By our first case above, we compare (i) and (ii), and (iii) and (iv) to see that $s=u=t$. Moreover, we have the same composition factors with each pair. But the latter is also true between (ii) and (iii), since the first $2$ composition factors $G/G_1,G_1/K_2$ or (ii), $H_0/H_1,H_1/K_2$ for (iii), which are isomorphic to each other.
\end{proof}

\begin{defn}
    A finite group $G$ is \textbf{solvable} if the simple factors in any composition series of $G$ are cyclic groups of prime order.
\end{defn}

\begin{ex}
    \begin{enumerate}
        \item Any abelian group is solvable
        \item $S_4$ is solvable, which has a composition series $1 \lhd C_2 \lhd V \lhd A_4 \lhd S_4$.
        \item $S_n$ is not solvable for $n \geq 5$, having composition series $1 \lhd A_n \lhd S_n$
    \end{enumerate}
\end{ex}

\begin{thm}[Feit-Thompson, 1963, 255 pages]
    If $|G|$ is odd, then $G$ is solvable. In particular, any non-abelian simple group must have even order.
\end{thm}

\begin{notebox*}[Question]
    If $N_1$ and $N_2$ are simple groups, how many groups $G$ have $N_1 \leq G$ and $G/N_1 \cong N_2$?
\end{notebox*}
This is too hard.

\section{Semidirect Products}

\begin{defn}[Semi-direct products]
    Suppose a group $G$ has a subgroup $H$ and a normal subgroup $N$, such that $G = NH$ and $N \cap H = \{1\}$. Recall that if $H$ were also normal, then $G$ would be the direct product of $N$ and $H$. Without this assumption, we say that $G$ is the \textbf{(internal) semi-direct product} $G = N \rtimes H$.
\end{defn}

\begin{rmk}
    \begin{enumerate}
        \item As $H \cap N=\{1\}$,
        $$H=H/H \cap N \cong HN /N = G/N,$$
        so if $G$ is finite, we have $|G| = |H| \cdot |N|$.
        \item Every $g \in G$ has a unique expression as $g=nh$ for $n \in N$ and $h \in H$.
        \item For $x,y \in G$, write $x=n_1h_1$ and $y=n_2h_2$, $n_i \in N,h_i \in H$. Then, $xy= n_1h_1n_2h_2=n_1(h_1n_2h_1^{-1})(h_1h_2)=n_1n_2'(h_1h_2)=nh$.
        \item For $h \in H$, we use conjugation by $h$ to define a map $\varphi_h:N \to N$ by $\varphi_h(n):=hnh^{-1},n \in N$. Then, $\varphi _h \in \Aut(N)$ and $\varphi_{h_1}\circ \varphi_{h_2}=\varphi_{h_1h_2}$, so we get a homomorphism $\varphi :H \to \Aut (N)$, $h \mapsto \varphi _h$. This $\varphi$ is called the \textbf{conjugation homomorphism} of the semi-direct product. We see that $(n_1h_1)(n_2h_2) = n_1\varphi(h_1)(n_2)\cdot h_1h_2$. So, the group operation on $G$ can be expressed in terms of the group operations on $N$ and $H$, and the homomorphism $\varphi :H \to \Aut (N)$
        \item If the conjugation homomorphism $\varphi$ is trivial, then
        $$nhn^{-1}=(nhn^{-1})(h^{-1}h)=n\varphi_h(n^{-1})h=n\varphi(h)(n^{-1})h=nn^{-1}h=h,$$
        so $H$ is normal in $G$, and $G = N \times H$. Conversely, if $G = N\times H$, then $\varphi$ is trivial. If $\varphi$ is non-trivial, then $G$ is non-abelian.
    \end{enumerate}
\end{rmk}

\section{External Semi Direct Products}

\begin{defn}
    Let $N$ and $H$ be groups, and $\varphi :H \to \Aut (N)$ any group homomorphism. Define a binary operation on $N \times H$ by $(n_1h_1)\cdot (n_2h_2) := (n_1\varphi(h_1)(n_2),h_1h_2)$. This turns $N \times H$ into a group with identity $(1,1)$; the inverse of $(n,h)$ is $(\varphi(h^{-1})(n^{-1}),h^{-1})$. This group is called the \textbf{external semi-direct product of $N$ by $H$, corresponding to $\varphi$}, denoted $G = N \rtimes _\varphi H$
\end{defn}

Now, $G$ has subgroups $\N := N \times 1$ and $\mathbb H = 1 \times H$, which are isomorphic to $N$ and $H$, respectively. For $(x,1) \in N$ and $(n,h) \in G$, we have
$$(n,h)(x,1)(n,h)^{-1}=(n\varphi(h)(x),h)(\varphi(h^{-1})(n^{-1}),h^{-1})=(n\varphi (h)(x) n^{-1},1) \in \N,$$
hence $\mathbb N \unlhd G$. Moreover, we see $(n,h) = (n\varphi(1)(1),h)=(n,1)\cdot (1,h)$, so $G = \mathbb N \cdot \mathbb H$. Finally, as $\N \cap \mathbb H = \{(1,1)\}$, we see that $G$ is the internal semi-direct product $G = \N \rtimes \mathbb H$.

Finally, given $(n,1) \in \N$ and $(1,h) \in \mathbb H$, we have $(1,h)(n,1)(1,h)^{-1}=(\varphi(h)(n),1)$, hence the conjugation homomorphism $\mathbb H \to \Aut (\mathbb N)$ of $\mathbb N \rtimes \mathbb H$ corresponds to the original homomorphism $\varphi :H \to \Aut (N)$. In this way, we can identify internal and external semi-direct products.

\begin{rmk}
    If $\varphi, \psi:H \to \Aut(N)$ are distinct homomorphisms, the groups $N \rtimes _\varphi H$ and $N \rtimes _\psi H$ are not necessarily isomorphic.
\end{rmk}

\begin{lem}[][label=lem:semidirect-iso]
    Let $H$ be a cyclic group, and $N$ be any group. Suppose that $\varphi ,\psi :H \to \Aut(N)$ are monomorphisms (1-1 homomorphisms) such that $\varphi (H) = \psi (H)$. Then, we have $N \rtimes _\varphi H \cong N \rtimes _\psi H$.
\end{lem}

\begin{proof}
    Let $H = \angles{x}$. Since $\varphi (H)$ and $\psi(H)$ coincide, we see that $\varphi(x)$ and $\psi(x)$ generate the same cyclic subgroup of $\Aut (N)$. Hence, there exists $a,b\in \Z$ such that $\varphi (x)^a = \psi(x)$, and $\psi(x)^b = \varphi(x)$. Since $H$ is cyclic generated by $x$, we get in general that $\varphi(h)^a=\psi (h),\psi(h)^b=\varphi(h)$ for all $h \in H$. Moreover, $\varphi(x)=\psi(x)^b=\varphi(x)^{ab}=\varphi(x^{ab})$, and $\varphi$ is 1-1, so $x^{ab}=x \iff h^{ab}=h$ for all $h \in H$. Now, define $\tau :N \rtimes _\psi H \to N \rtimes _\varphi H$ by $\tau (nh) := nh^a$.

    Then,
    $$\tau((n_1h_1)(n_2h_2))=\tau (n,\psi(h_1)(n_2)h_1h_2)=\cdots =(n_1h_1^a)(n_2h_2^a)=\tau(n_1h_1)\tau(n_2h_2)$$
    Similarly, we show that $\tau ' : N \rtimes _\varphi H \to N \rtimes_\psi H$ sending $nh$ to $nh^b$ is a homomorphism. Since $h^{ab}=h$, it follows that $\tau$ and $\tau '$ are mutually inverse, hence they are isomorphisms.
\end{proof}

\begin{prop}
    If $G = C_n$ is a cyclic group of order $n$, then $\Aut G \cong (\Z/n\Z)^\times$, an abelian group of order $\varphi(n)$.
\end{prop}

\begin{proof}
    Let $g$ be a generator of $G=C_n$. If $\varphi \in \Aut G$, then $\varphi(g)=g^a$ uniquely determines $\varphi$; we denote this by $\varphi_a$. Since $o(g^a)=o(g)=n$, we must have $(a,n) = 1$, and conversely for any such $a$ coprime to $n$, the map $\varphi(x)=x^a$ gives an automorphism of $C_n$. We thus have a surjective map $f:\Aut (C_n) \to (\Z/n\Z)^\times$ mapping $\varphi_a \mapsto a+n\Z$, and $f$ is a group homomorphism. Hence, it is an isomorphism. We'll see later than if $n=p$ is prime, then $\Aut (C_p) \cong C_{p-1}$. More generally, this remains true for $\Aut(C_{p^k})\cong C_{p^{k-1}(p-1)}$ where $p \geq 3$. However, for all $k \geq 3$, $\Aut(C_{2^k}) \cong C_2 \times C_{2^{k-2}}$.
\end{proof}

\section{\texorpdfstring{Groups of Order $pq$}{Groups of Order pq}}

Let $p,q$ be distinct primes, $p>q$.

\begin{thm}
    If $p \not \equiv 1 \pmod q$, then any group of order $pq$ is isomorphic to $C_{pq}$. Otherwise, if $p \equiv 1 \pmod q$, then any abelian group of order $pq$ is isomorphic to $C_{pq}$, and there's exactly one non-abelian group of order $pq$, up to isomorphism.
\end{thm}

\begin{proof}
    Let $|G|=pq$, and let $P$ and $Q$ be a $p$-Sylow and $q$-Sylow subgroup of $G$, respectively. Since $|P| = p,|Q|=q$, we have $P \cong C_p,Q \cong C_q$, and $P \cap Q= \{1\}$. Now, the number of conjugates of $P$ in $G$ divides $|G:P|=q$, and is congruent to $1$ modulo $p$. But, $q \not \equiv 1 \pmod p$ as $q<p$, so $P$ has only one conjugate in $G$, hence $P \lhd G$.

    Now, $Q$ has either $1$ or $p$ conjugates in $G$, and if the latter happens, we must have $p \equiv 1 \pmod q$. If $Q \unlhd G$, then $G= C_p \times C_q \cong C_{pq}$. So, we are reduced to the case when $p \equiv 1 \pmod q$ and $G$ is non-abelian. In this instance, $G = P \rtimes Q$. Let $\varphi : Q \to \Aut (P)$ be the conjugation homomorphism. Now, $\varphi$ is non-trivial, and $Q$ is a simple group, so $\varphi$ is $1-1$. We're reduced to studying monomorphisms  $C_q \to \Aut( C_p)=C_{p-1}$. Since $q \mid p-1$, $\Aut(C_p)$ has a unique subgroup $K$ of order $q$. Define $\varphi:C_q \to \Aut (C_p)$ by letting $\varphi(x)$ be a generator of the cyclic group $K \leq \Aut (C_p)$, where $x$ is a generator of $C_q$. This shows that a nontrivial $\varphi$ exists, so a non-abelian group of order $pq$.

    Now, if $\psi :C_q \to \Aut(C_p)$ is another monomorphism, then by uniqueness of $K$, we must have $\psi(C_q) = K = \varphi (C_q)$. Hence, \Cref{lem:semidirect-iso} shows that $C_p \rtimes_\varphi C_q \cong C_p \rtimes _\psi C_q$, this is non-abelian group is unique.
\end{proof}

\begin{ex}[Examples of Semi-Direct Products]
    \begin{enumerate}
        \item $M_n$ is the group of rigid motions of $\R^n$. Then, $M_n = O(n) \ltimes \R^n$.
        \item The dihedral group $D_n = C_n \rtimes C_2$
        \item The symmetric group $S_n=A_n \rtimes C_2$.
    \end{enumerate}
\end{ex}

\section{The Free Group}

\begin{notebox*}[Motivation]
    Make precise presentations of groups like
    $$C_n=\angles{x \mid x^n=1} \quad \text{or} \quad D_n = \angles{x,y\mid x^2=y^n=1,xy=y^{-1}x}.$$
\end{notebox*}

\begin{defn}[Words]
    Let $S$ be a (possibly infinite) set of symbols (e.g. $\{a,b,\dots\}$), and define $S'=S\cup S^{-1}=\{a,a^{-1},b,b^{-1},\dots\}$. A \textbf{word} is a finite string of symbols of $S'$. Words are multiplied by concatenation:
    $$(aa)(ba)=aaba$$
    $1$ denotes the \textbf{empty word} with no letters, and let $W$ be the set of all words.

    If a word contains $xx^{-1}$ or $x^{-1}x$ for some $x \in S$, we agree to \textbf{cancel} $x,x^{-1}$, and reduce the length of the word. The word is \textbf{reduced} if no such cancellation can occur. Starting with any $w \in W$, we can perform a sequence of cancellations until we reach a reduced word $w_0$. This is called a \textbf{reduced form} of $w$.
\end{defn}

\begin{prop}
    There is only one reduced form for any given $w \in W$
\end{prop}

\begin{proof}
    We may induct on the length of $w$. If $w$ is already reduced, then there's nothing to prove. Otherwise, there exists a pair $xx^{-1}$ in $W$ which can be canceled (here, $x \in S'$, and by convention if $x = a^{-1}$, then $x^{-1}=a$). It suffices, using induction, to prove that we can obtain each reduced form of $w$ by canceling $xx^{-1}$ first.

    Let $w_0$ be a reduced form of $w$, i.e. $w_0$ is obtained from $w$ by a sequence of cancellations. If $xx^{-1}$ is canceled at some point during the process, we can rearrange to cancel it first. The only other possibility is that $x$ or $x^{-1}$ in "$xx^{-1}$" was canceled at some point. If, say, $x$ was canceled, then it happened like $x^{-1}xx^{-1} \to x^{-1}$. But, notice that canceling the original pair has the same effect. (similar if $x^{-1}$ was canceled first) So, we're back in the previous case.
\end{proof}

\begin{defn}
    $w,w'\in W$ are \textbf{equivalent} if the have the same reduced form.
\end{defn}

\begin{prop}
    The product of equivalent words are equivalent. We can therefore multiply equivalence classes of words.
\end{prop}

\begin{proof}
    If $w \sim w',v \sim v'$, let $w_0,v_0$ be the reduced form for $w,w'$, and $v,v'$ respectively. It follows that $wv \sim w_0v_0 \sim w'v'$.
\end{proof}

\begin{rmk}
    We have inverses, as, for example $[abc\cdots z]^{-1}=[z^{-1}\cdots c^{-1}b^{-1}a^{-1}]$
\end{rmk}

\begin{defn}[The Free Group]
    The group $F(S) = \angles{S}$ of equivalence classes of words on $S'$ is the \textbf{free group} on $S$.
\end{defn}

We can now introduce power notation:
$$aaabb=a^3b^2.$$
\begin{defn}
    The elements of $S$ are the \textbf{generators} of $F(S)$
\end{defn}

\begin{ex}
    \begin{enumerate}
        \item If $S= \{x\}$, then $F(S) = \{x^k \mid k \in \Z\}$ is the infinite cyclic group generated by $x$.
        \item If $S = \{a,b\}$, $F(a,b) = \{1,a,a^{-1},b,b^{-1},ab,ba,a^{-1}b,\dots\}$
    \end{enumerate}
\end{ex}

\begin{thm}[Universal Property of the Free Groups]
    Notice that the free group $F(S)$ comes with a map of sets $f:S \to F(S)$ with $a \mapsto a$. The pair $(S,F(S))$ has the following universal property:

    For any group $G$ and map of sets $\varphi:S \to G$, there exists a unique homomorphism $\tilde \varphi :F(S) \to G$ such that $\tilde \varphi \circ f = \varphi$. That is, the below diagram commutes.
\end{thm}

% https://q.uiver.app/#q=WzAsMyxbMCwwLCJTIl0sWzIsMCwiRihTKSJdLFsxLDIsIkciXSxbMCwyLCJcXHZhcnBoaSIsMl0sWzAsMSwiZiJdLFsxLDIsIlxcZXhpc3RzIFxcdGlsZGUgXFx2YXJwaGkiLDAseyJzdHlsZSI6eyJib2R5Ijp7Im5hbWUiOiJkYXNoZWQifX19XV0=
\[\begin{tikzcd}
	S && {F(S)} \\
	\\
	& G
	\arrow["f", from=1-1, to=1-3]
	\arrow["\varphi"', from=1-1, to=3-2]
	\arrow["{\exists \tilde \varphi}", dashed, from=1-3, to=3-2]
\end{tikzcd}\]

\begin{proof}
    Not the whole proof, but observe that the obvious definition of $\tilde \varphi$ on words in $W$ passes to equivalence classes.
\end{proof}

\begin{rmk}
    If $G$ is a group and $S \subset G$ is any subset, from the universal property we get a unique map $\varphi:F(S) \to G$. $\varphi$ is surjective if and only if $G$ is generated by $S$. When $\varphi$ is onto, we get $G \cong F(S) / \ker \varphi$, hence the elements in $\ker \varphi$ are relations among the generators in the set $S$.
\end{rmk}

\begin{defn}[Relations]
    A set of words $R=\{r_1,r_2,\dots\}$ is a set of \textbf{defining relations} if $R \subset \ker \varphi$ from the above remark, and $\ker \varphi$ is the smallest normal subgroup containing $R$ (so, $\ker\varphi$ is generated by all words in $R$ with their conjugates).

    We use the notation $G= \angles{x_1,x_2, \dots  \mid r_1,r_2,\dots}$, where $S = \{x_1,x_2,\dots\}$.
\end{defn}

\begin{ex}
    \begin{enumerate}
        \item $G = \Z/n\Z$, $\varphi :\angles{x} \to G$, $x \mapsto 1$. $\ker \varphi$ is generated by $x^n$, so $G = \angles{x \mid x^n=1}$
        \item $G = \angles{x,y \mid xyx^{-1}y^{-1} = 1}= \{x^iy^j \mid i,j \in \Z\} \cong \Z \oplus \Z$, the \textbf{free abelian group} on $x$ and $y$.
    \end{enumerate}
\end{ex}

\textbf{This is where the material for the midterm ends.}

\section{Reflection Groups and Coxeter Groups}
\begin{defn}[Reflections]
    Let $V = \R^n$. For any $\alpha \neq 0$ in $V$, let $H_\alpha = \angles{\alpha}^{\perp}=\{ x \in V \mid \angles{x ,\alpha}=0\}$. Then, the \textbf{reflection} $s_\alpha:V \to V$ is the map in $O(V)\cong O(n)$ such that $s_\alpha (x)=x$ for all $x \in H_\alpha$, and $s_\alpha (\alpha) = -\alpha$.
\end{defn}

\begin{prop}
    $s_\alpha (x) = x- \frac{2\angles{x,\alpha}}{\angles{\alpha,\alpha}}\alpha$
\end{prop}

\begin{defn}[The Reflection Group]
    The reflection group is a finite subgroup $W$ of $O(V)$ generated as a group of reflections.
\end{defn}

\begin{ex}
    \begin{enumerate}
        \item $D_n = \angles{x,y \mid x^2=1,y^n=1,(xy^2)=1}$. So, set $s_1 = x,s_2=xy$, and $D_n = \angles{s_1,s_2 \mid s_1^2=s_2^2=1,(s_1s_2)^n=1}$
        \item The symmetric group $S_n \leq O(n)$, where $w \in S_n$ permutes $e_1,\dots,e_n$ ($w$ is a permutation matrix). The transposition $(ij)$ acts as the reflections $s_{e_i-e_j}$ about $H_{e_i-e_j}=\{x \in \R^n \mid x_i=x_j\}$. In fact, $S_n$ is the symmetry group of the $(n-1)$-simplex $\Delta_{n-1}:=\{\sum_{i=1}^n\lambda_ie_j \mid \lambda_i \geq 0, \lambda_1+\cdots + \lambda_n=1\}$.
        $$S_n = \{s_1,\dots, s_{n-1} \mid (s_is_j)^2=1,|i-j|\geq 2,(s_is_{i+1})^3=1\}$$
        \item Let $V = \R^n$, and let $S_n$ act as before. Now, add the $n$ reflections about the coordinate hyperplanes $H_{e_i}$, $1 \leq i \leq n$. These $n$ reflections generate a group $N$ of order $2^n$ isomorphic to $(\Z/2\Z)^n$. Note that $N \cap S_n=1$, and $N$ and $S_n$ generate a group $W=N \rtimes S_n$, called the \textbf{hyperoctahderal group} (the permutation action of $S_n$ on $(\Z/2\Z)^n$ is the conjugation homomorphism $S_n \to \Aut (N)$). $W$ is the symmetry group of the $n$-cube $[-1,1]^n=\{ \sum \lambda _ie_i \mid -1 \leq \lambda _i \leq 1\}$. The elements of $W_n$ are \textbf{signed permutation matrices} (like permutation matrices but they can have $\pm 1$ instead of just $1$). In one line notation, $w \in W$ is a \textbf{signed permutation} $w=(w_1,w_2,\dots,w_n)$, where $w_i = w(i)$, and $\{|w_1|,\dots, |w_n|\}=\{1,\dots,n\}$. $W_n$ is generated by the simple reflections $s_1,\dots,s_{n-1} \in S_n$ and the \textbf{sign change} $s_0:= (-1,2,\dots,n)$. The presentation is $W_n = \{s_0,\dots ,s_{n-1} \mid s_i^2=1,(s_is_j)^2=1,|i-j|\geq 2,(s_is_{i+1})^3=1,i\geq 1,(s_0s_1)^4=1\}$
    \end{enumerate}
\end{ex}

\begin{lem}
    Given $\alpha,\beta \in \R^2 \setminus \{0 \}$, $s_\beta \circ s_\alpha$ is a rotation through twice the angle between $\alpha$ and $\beta$. There's a version of this in $\R^n$ as well.
\end{lem}
\begin{proof}
    Exercise
\end{proof}

\begin{defn}[Coxeter Groups]
    A group $W$ is a \textbf{Coxeter Group} if there exists $S \subset W$ such that $W$ is a quotient of the free group $F(S)$ modulo the relations:
    \begin{enumerate}[(i)]
        \item $s^2=1$ for all $s \in S$,
        \item $(ss')^{m_{s,s'}}=1$ for all $s,s' \in S$ with $ s \neq s'$, and for some minimal integers $m_{s,s'}=m_{s',s}\geq 2$.
    \end{enumerate}
    We represent $W$ by a labeled graph $\Gamma _W$ called the \textbf{Coxeter graph} with one vertex for each $s \in S$, and an edge connecting $s,s'$ if $m_{s,s'}\geq3$ labeled by $m_{s,s'}$. In practice, if $m_{s,s'}=3$, the label $3$ is omitted.

    $W$ is \textbf{irreducible} if $\Gamma_W$ is \textbf{connected} (if not, we can write $W = W_1\times \cdots \times W_r$ depending on the connected components of $\Gamma _W$).
\end{defn}

\begin{ex}
    \begin{enumerate}
        \item $D_{m}=I_2(m)$
        \item $S_n \equiv A_{n-1}$
        \item $W_n \equiv BC_n$
    \end{enumerate}
\end{ex}

\begin{thm}
    Let $(W,S)$ be a Coxeter group, where $S$ is a finite set of Coxeter generators, with connected Coxeter diagram $\Gamma_W$. Then, TFAE
    \begin{enumerate}
        \item $W$ is a finite group
        \item $\Gamma$ is one of the following Coxeter diagrams: $A_n$ with $n\geq 1$ vertices, $BC_n$ with $n \geq 2$ vertices, $D_n$ with $n \geq 4$ vertices, $E_6$, $E_7$, $E_8$, $F_4$, $G_2$, $H_3$, $H_4$, $I_2(m)$. See below.
    \end{enumerate}
\end{thm}

\begin{figure}[H]
    \centering
    \includegraphics[width=\linewidth]{coxeter-diagrams.png}
    \caption{Source: Rgugliel on \href{https://en.wikipedia.org/wiki/Coxeter_group}{Wikipedia}}
\end{figure}
